% Chapter 1 of the BARCS student guide: the ideas.
%=====================================================================
\section{The experiment: a CRISPR screen}
%=====================================================================

Each cell in your body has about 20{,}000 genes. A gene is a set of
instructions. Scientists want to know which genes a cell needs, for example
to survive a cancer drug.

A \textbf{CRISPR screen} tests thousands of genes at the same time. It works
in four steps:

\begin{enumerate}
\item Scientists make a \textbf{library} of short pieces of RNA, called
  \textbf{guides}. Each guide switches off one gene.
\item They put one guide into each cell. Now each cell has one gene switched
  off.
\item They grow the cells. Some cells grow under normal conditions
  (\textbf{control}). Other cells grow with a drug (\textbf{drug}).
\item A sequencing machine reads the guides in each group. The machine
  counts how many times it reads each guide. This number is a
  \textbf{read count}.
\end{enumerate}

If cells need a gene to survive the drug, the cells that lose that gene die.
Their guide becomes rare in the drug group. So the guide count goes down.

\begin{remember}
BARCS does not look at cells directly. It looks at read counts. A guide that
becomes rare tells us that its gene is important.
\end{remember}

Scientists repeat each group more than one time. Each repeat is a
\textbf{replicate}. Each sample that the machine reads is a \textbf{library}.
In this guide, we use four libraries: two control replicates and two drug
replicates.

%=====================================================================
\section{The data}
%=====================================================================

The machine gives two numbers for each library:

\begin{itemize}
\item the read count of our guide, and
\item the total read count of all guides in that library.
\end{itemize}

Table~\ref{tab:data} shows the data for one guide. We use these numbers in
all sections of this guide.

\begin{table}[h]
\centering
\caption{Read counts for one guide in four libraries.}
\label{tab:data}
\begin{tabular}{lcrr}
\toprule
Library & Group & Guide count $K$ & Total count $N$\\
\midrule
1 & control & 180 & 1{,}000{,}000\\
2 & control & 240 & 1{,}200{,}000\\
3 & drug    & 120 & 1{,}000{,}000\\
4 & drug    &  80 &   800{,}000\\
\bottomrule
\end{tabular}
\end{table}

\begin{symbols}
\begin{tabular}{@{}l p{0.78\linewidth}@{}}
$K$ & The read count of one guide in one library. Library~1 has $K=180$.\\
$N$ & The total read count of all guides in one library. Library~1 has
      $N=1{,}000{,}000$.\\
$i$ & The library number (1, 2, 3 or 4). We write $K_i$ for the count in
      library $i$. So $K_2=240$.\\
\end{tabular}
\end{symbols}

\begin{remember}
Always record $N$ before you remove any guides from the table. If you remove
guides first, the totals become wrong. Then all the later numbers become
wrong too.
\end{remember}

%=====================================================================
\section{Step 1: Turn counts into proportions}
%=====================================================================

Library~2 has more reads than library~1 because the machine read it more.
So we cannot compare raw counts. We divide each count by its library total.
The result is a \textbf{proportion}:
\[
p = \frac{K}{N}.
\]

Proportions are very small numbers. So we multiply by one million. The
result is the count \textbf{per million} reads.

\begin{worked}[ 1: proportions]
\begin{tabular}{@{}lll@{}}
Library 1: & $p = 180 \div 1{,}000{,}000 = 0.000180$ & $= 180$ per million\\
Library 2: & $p = 240 \div 1{,}200{,}000 = 0.000200$ & $= 200$ per million\\
Library 3: & $p = 120 \div 1{,}000{,}000 = 0.000120$ & $= 120$ per million\\
Library 4: & $p = \phantom{0}80 \div \phantom{0,}800{,}000 = 0.000100$ & $= 100$ per million\\
\end{tabular}

\smallskip
The two control libraries give 180 and 200 per million. The two drug
libraries give 120 and 100 per million. The guide is less common with the
drug.
\end{worked}

\begin{symbols}
\begin{tabular}{@{}l p{0.78\linewidth}@{}}
$p$ & The proportion of all reads that belong to the guide. It is always
      between 0 and 1.\\
\end{tabular}
\end{symbols}

%=====================================================================
\section{Step 2: Compare the groups with a fold change}
%=====================================================================

We want one number that tells how much the guide changes. First, we combine
the replicates of each group. We add their counts and add their totals:

\[
p_{\text{control}} = \frac{180+240}{1{,}000{,}000+1{,}200{,}000}
 = \frac{420}{2{,}200{,}000} = 0.000191 \;(191 \text{ per million}),
\]
\[
p_{\text{drug}} = \frac{120+80}{1{,}000{,}000+800{,}000}
 = \frac{200}{1{,}800{,}000} = 0.000111 \;(111 \text{ per million}).
\]

The \textbf{fold change} is the drug proportion divided by the control
proportion:
\[
\text{fold change} = \frac{p_{\text{drug}}}{p_{\text{control}}}
 = \frac{0.000111}{0.000191} = 0.58.
\]

A fold change of 0.58 means that the guide keeps only 58\% of its control
level. It loses 42\%.

\begin{remember}
\begin{itemize}[leftmargin=1.2em]
\item Fold change $= 1$: no change.
\item Fold change $< 1$: the guide becomes rarer. Cells may need this gene.
\item Fold change $> 1$: the guide becomes more common. Losing this gene
  may help the cells.
\end{itemize}
\end{remember}

\subsection{Logarithms make changes symmetric}

A fold change of 2 means ``double''. A fold change of 0.5 means ``half''.
These changes have the same size, but the numbers 2 and 0.5 do not look
equal. A \textbf{logarithm} fixes this problem.

The \textbf{base-2 logarithm}, $\log_2$, tells how many times you double.
\begin{center}
\begin{tabular}{lccccc}
\toprule
Fold change & 0.25 & 0.5 & 1 & 2 & 4\\
$\log_2$(fold change) & $-2$ & $-1$ & 0 & 1 & 2\\
\bottomrule
\end{tabular}
\end{center}
Now ``double'' is $+1$ and ``half'' is $-1$. The sizes look equal.

\begin{worked}[ 2: log fold change]
Our fold change is 0.58. It is between 0.5 and 1. So $\log_2(0.58)$ is
between $-1$ and 0. A calculator gives $\log_2(0.58) = -0.78$.

The guide drops by 0.78 ``halvings''.
\end{worked}

BARCS uses the \textbf{natural logarithm}, $\ln$. It counts powers of the
number $e = 2.718\ldots$ instead of powers of 2. The idea is the same. To
change back from $\ln$, you raise $e$ to that power: if $\ln(x) = b$, then
$x = e^{b}$. For our guide, $\ln(0.58) = -0.54$.

%=====================================================================
\section{Step 3: Odds and log-odds}
%=====================================================================

BARCS works with \textbf{odds} instead of proportions. The odds compare the
reads that belong to the guide with the reads that do not:
\[
\text{odds} = \frac{p}{1-p}.
\]
The \textbf{log-odds} is the natural logarithm of the odds:
\[
\text{log-odds} = \ln\!\left(\frac{p}{1-p}\right).
\]

\begin{worked}[ 3: odds and log-odds]
Control: $p = 0.000191$. So
$\text{odds} = 0.000191 \div 0.999809 = 0.000191$.

The odds and the proportion are almost equal. This is true for all small
proportions, because $1-p$ is almost 1.

Log-odds for control: $\ln(0.000191) = -8.56$.\\
Log-odds for drug: $\ln(0.000111) = -9.10$.\\
Difference: $-9.10 - (-8.56) = -0.54$.

The difference of the log-odds is $-0.54$. This is the same as
$\ln(\text{fold change})$ from Step~2.
\end{worked}

\begin{remember}
For guides, a difference in log-odds is almost exactly the natural log of
the fold change. So you can read BARCS results as log fold changes.
\end{remember}

Why use log-odds at all? A proportion must stay between 0 and 1. A log-odds
can be any number, positive or negative. This property lets BARCS add
effects together safely, as Step~5 shows.

%=====================================================================
\section{Step 4: Two kinds of noise}
%=====================================================================

The two control libraries gave 180 and 200 per million. They are not equal.
Replicates always differ a little. We call this difference
\textbf{noise}. BARCS must know how big the noise is. Then it can tell a
real change from a chance change.

Guide counts have two kinds of noise:

\begin{enumerate}
\item \textbf{Counting noise.} The machine reads a random sample of all
  the guides. Think of a bag with a million marbles. You take out a handful.
  You do not get exactly the same number of red marbles each time.
\item \textbf{Biological noise.} Each replicate is a different group of
  living cells. The cells grow a little differently each time.
\end{enumerate}

BARCS measures the size of the noise with the \textbf{variance}. The
variance is the average squared distance of a value from its expected
value. The \textbf{standard deviation} is the square root of the variance.
It is the typical distance in the same units as the counts.

BARCS uses this formula for the variance of a guide count:
\[
\text{variance} = \underbrace{N\,\mu\,(1-\mu)}_{\text{counting noise}}
 \times \underbrace{\bigl(1 + (N-1)\,\rho\bigr)}_{\text{extra biological noise}}.
\]

\begin{symbols}
\begin{tabular}{@{}l p{0.78\linewidth}@{}}
$\mu$ & ``mu''. The true proportion of the guide that we expect. For the
        control group, our best value is $\mu = 0.000191$.\\
$N$ & The library total, as before.\\
$\rho$ & ``rho''. A number that measures the biological noise. If
         $\rho = 0$, only counting noise exists. A bigger $\rho$ means more
         biological noise.\\
\end{tabular}
\end{symbols}

\begin{worked}[ 4: variance]
Use $N = 1{,}000{,}000$ and $\mu = 0.000191$. For our guide, BARCS
estimates $\rho = 0.000000386$.

\textbf{Counting noise:}
$N\mu(1-\mu) = 1{,}000{,}000 \times 0.000191 \times 0.999809 = 191$.

\textbf{Extra factor:}
$1 + (N-1)\rho = 1 + 999{,}999 \times 0.000000386 = 1.39$.

\textbf{Total variance:} $191 \times 1.39 = 265$.

\textbf{Standard deviation:} $\sqrt{265} = 16.3$ reads.

So a count of about 191 normally moves up or down by about 16. Biological
noise adds 39\% to the counting noise.
\end{worked}

\begin{remember}
Look at the extra factor $1 + (N-1)\rho$. It contains $N$. When the machine
reads more, $N$ grows, and the factor grows too. So more reading does not
remove biological noise. Only more replicates help with biological noise.
\end{remember}

%=====================================================================
\section{Step 5: The design matrix}
%=====================================================================

A \textbf{model} is a rule that predicts a value. BARCS predicts the
log-odds of a guide in each library. The model needs to know which group
each library belongs to. BARCS gets this information from the
\textbf{design matrix}.

\subsection{What a design matrix is}

A design matrix is a table of numbers. It has:
\begin{itemize}
\item one \textbf{row} for each library, and
\item one \textbf{column} for each $\beta$ value in the model.
\end{itemize}
We write the design matrix with the letter $X$.

For our four libraries, the model has two $\beta$ values. So $X$ has four
rows and two columns:

\begin{center}
\begin{tabular}{lc|cc}
\toprule
Library & Group & Column 1: ``intercept'' & Column 2: ``drug''\\
\midrule
1 & control & 1 & 0\\
2 & control & 1 & 0\\
3 & drug    & 1 & 1\\
4 & drug    & 1 & 1\\
\bottomrule
\end{tabular}
\end{center}

\begin{symbols}
\begin{tabular}{@{}l p{0.74\linewidth}@{}}
$X$ & The design matrix. One row for each library, one column for each
      $\beta$.\\
Intercept & A column of 1s. It belongs to $\beta_0$, the starting level.
      Every library gets the starting level.\\
Drug & A column that is 1 for a drug library and 0 for a control library.
      It belongs to $\beta_1$, the effect of the drug.\\
$\beta_0$ & ``beta zero''. The log-odds of the guide in the control group.\\
$\beta_1$ & ``beta one''. The change in log-odds from control to drug. This
      is the \textbf{effect}. It is the most important number.\\
\end{tabular}
\end{symbols}

\subsection{From the design matrix to a prediction}

To get the prediction for one library, take the row of that library. Multiply
each number in the row by its $\beta$. Then add the products:
\[
\text{predicted log-odds} = (\text{column 1}) \times \beta_0 +
(\text{column 2}) \times \beta_1 .
\]
This calculation is a \textbf{matrix multiplication}. You do it one row at
a time. We call the predicted log-odds $\eta$ (``eta'').

\begin{worked}[ 5: predictions from the design matrix]
BARCS finds $\beta_0 = -8.565$ and $\beta_1 = -0.542$ (Step~6 in Chapter~2 shows how).

\begin{tabular}{@{}lll@{}}
Library 1 (row 1, 0): & $1 \times (-8.565) + 0 \times (-0.542)$ & $= -8.565$\\
Library 2 (row 1, 0): & $1 \times (-8.565) + 0 \times (-0.542)$ & $= -8.565$\\
Library 3 (row 1, 1): & $1 \times (-8.565) + 1 \times (-0.542)$ & $= -9.107$\\
Library 4 (row 1, 1): & $1 \times (-8.565) + 1 \times (-0.542)$ & $= -9.107$\\
\end{tabular}

\smallskip
The control libraries get $\beta_0$ only. The drug libraries get
$\beta_0 + \beta_1$. So $\beta_1$ is the difference between the groups.
\end{worked}

\subsection{From log-odds back to counts}

A prediction in log-odds is hard to read. So we change it back into a
proportion, and then into a count. The rule is:
\[
\mu = \frac{e^{\eta}}{1 + e^{\eta}}, \qquad
\text{expected count} = N \times \mu .
\]
This rule undoes the log-odds of Step~3.

\begin{worked}[ 6: expected counts]
\textbf{Control libraries} ($\eta = -8.565$):
$e^{-8.565} = 0.000190676$. Then
$\mu = 0.000190676 \div 1.000190676 = 0.000190640$, which is 190.6 per
million.

\textbf{Drug libraries} ($\eta = -9.107$):
$e^{-9.107} = 0.000110840$. Then $\mu = 0.000110827$, which is 110.8 per
million.

\begin{tabular}{@{}lrrrr@{}}
\toprule
Library & $N$ & $\mu$ (per million) & Expected count $N\mu$ & Real count $K$\\
\midrule
1 & 1{,}000{,}000 & 190.6 & 190.6 & 180\\
2 & 1{,}200{,}000 & 190.6 & 228.8 & 240\\
3 & 1{,}000{,}000 & 110.8 & 110.8 & 120\\
4 &   800{,}000 & 110.8 &  88.7 &  80\\
\bottomrule
\end{tabular}

\smallskip
The expected counts are near the real counts. The differences are
$-10.6$, $+11.2$, $+9.2$ and $-8.7$ reads. These differences are the noise.
\end{worked}

\subsection{The reference group}

The control group has 0 in the drug column. So control is the
\textbf{reference group}. The effect $\beta_1$ always compares a group with
the reference group. If you make drug the reference group, $\beta_1$ changes
sign. It becomes $+0.542$. The meaning stays the same.

\subsection{Paired replicates: add a column}

Sometimes the replicates are pairs. For example, replicate~1 of control and
replicate~1 of drug come from the same batch of cells. Then BARCS can remove
the batch difference. You add a column for the batch:

\begin{center}
\begin{tabular}{lcc|ccc}
\toprule
Library & Group & Replicate & Intercept & Replicate 2 & Drug\\
\midrule
1 & control & 1 & 1 & 0 & 0\\
2 & control & 2 & 1 & 1 & 0\\
3 & drug    & 1 & 1 & 0 & 1\\
4 & drug    & 2 & 1 & 1 & 1\\
\bottomrule
\end{tabular}
\end{center}

Now the model has three $\beta$ values: the starting level, the batch
difference and the drug effect. Each new column costs one degree of freedom
(Step~8 in Chapter~2). With four libraries and three columns, only one degree of freedom
stays.

\begin{worked}[ 7: the paired model for our guide]
BARCS gives the drug effect $\beta = -0.544$ with standard error 0.143.
The effect is almost the same as before ($-0.542$). But the degrees of
freedom drop from 2 to 1. The p-value becomes 0.164 instead of 0.033.

Our two replicates have almost no batch difference ($\beta = 0.005$). So
the extra column costs information and gives nothing back. Add a batch
column only when the replicates really are pairs.
\end{worked}

\subsection{Numbers instead of groups}

A column can also hold a number, for example the day of the experiment.
Suppose that you collect cells on day 0, day 7 and day 14, two libraries
each. You put the day divided by 14 into a column:

\begin{center}
\begin{tabular}{lc|cc}
\toprule
Library & Day & Intercept & Day $\div$ 14\\
\midrule
1, 2 & 0  & 1 & 0\\
3, 4 & 7  & 1 & 0.5\\
5, 6 & 14 & 1 & 1\\
\bottomrule
\end{tabular}
\end{center}

Now $\beta_1$ is the change in log-odds over 14 days. The model draws a
straight line through the three days.

\begin{worked}[ 8: a time course]
Suppose that $\beta_1 = -0.70$.

After 14 days: change $= 1 \times (-0.70) = -0.70$. Fold change
$= e^{-0.70} = 0.50$. The guide drops to half.

After 7 days: change $= 0.5 \times (-0.70) = -0.35$. Fold change
$= e^{-0.35} = 0.70$. The guide drops to 70\%.

The model has 6 libraries and 2 columns. So it has $6 - 2 = 4$ degrees of
freedom.
\end{worked}

\subsection{Two factors and an interaction}

Some experiments change two things. For example, cells can be normal or can
have a second gene switched off (``knockout''). And cells can get no drug or
the drug. This experiment has four groups. The design matrix has four
columns:

\begin{center}
\begin{tabular}{l|cccc}
\toprule
Group & Intercept & Knockout & Drug & Knockout $\times$ Drug\\
\midrule
normal, no drug   & 1 & 0 & 0 & 0\\
knockout, no drug & 1 & 1 & 0 & 0\\
normal, drug      & 1 & 0 & 1 & 0\\
knockout, drug    & 1 & 1 & 1 & 1\\
\bottomrule
\end{tabular}
\end{center}

The last column is the product of the knockout column and the drug column.
It is 1 only when both are 1. Its $\beta$ is the \textbf{interaction}. The
interaction measures the extra change that happens only when both changes
happen together.

\begin{worked}[ 9: reading an interaction]
Suppose that the four $\beta$ values are $-8.5$, $-0.2$, $+0.1$ and
$-0.7$. Multiply each row by the $\beta$ values and add:

\begin{tabular}{@{}ll@{}}
normal, no drug: & $-8.5$\\
knockout, no drug: & $-8.5 + (-0.2) = -8.7$\\
normal, drug: & $-8.5 + 0.1 = -8.4$\\
knockout, drug: & $-8.5 + (-0.2) + 0.1 + (-0.7) = -9.3$\\
\end{tabular}

\smallskip
The knockout alone changes the log-odds by $-0.2$. The drug alone changes it
by $+0.1$. If the two effects only added up, the last group would be
$-8.5 - 0.2 + 0.1 = -8.6$. The real value is $-9.3$. The extra $-0.7$ is the
interaction. The guide drops strongly only when both changes happen.
\end{worked}

\subsection{Comparing two treatments: a contrast}

Suppose an experiment has three groups: control, drug A and drug B. The
design matrix has the columns intercept, drug A and drug B. Then $\beta_A$
compares drug A with control, and $\beta_B$ compares drug B with control. To
compare drug A with drug B, you subtract: $\beta_A - \beta_B$. This
calculation is a \textbf{contrast}. In BARCS, the function
\texttt{bb\_contrast()} calculates it and its test.

For example, if $\beta_A = -0.9$ and $\beta_B = -0.4$, then
$\beta_A - \beta_B = -0.5$. Drug A lowers the guide more than drug B.

\subsection{Rules for a good design matrix}

\begin{enumerate}
\item Make one row for each library, in the same order as the count table.
\item Make one column for each $\beta$ that you want.
\item Keep more rows than columns. The difference is the degrees of
  freedom. Two or more is better than one.
\item Do not let two columns carry the same information. For example,
  suppose that all drug libraries come from batch~2. Suppose also that all
  control libraries come from batch~1. Then the batch column and the drug
  column are equal. BARCS cannot separate the drug effect from the batch
  effect.
\end{enumerate}

\subsection{The design matrix in R}

In R, you describe the design matrix with a \textbf{formula}. The formula
\texttt{\textasciitilde\ group} means ``an intercept and a column for each
group except the reference group''.

The function \texttt{model.matrix()} builds $X$ from the formula:

\begin{center}
\begin{minipage}{0.86\textwidth}
\small
\begin{verbatim}
> data <- data.frame(group = c("control", "control", "drug", "drug"))
> model.matrix(~ group, data)
  (Intercept) groupdrug
1           1         0
2           1         0
3           1         1
4           1         1
\end{verbatim}
\end{minipage}
\end{center}

Then you give the counts, the totals, the formula and the column name to
BARCS:

\begin{center}
\begin{minipage}{0.86\textwidth}
\small
\begin{verbatim}
fit <- bb_screen(counts, data, formula = ~ group,
                 term = "groupdrug", totals = totals)
\end{verbatim}
\end{minipage}
\end{center}

For the paired design, use the formula \texttt{\textasciitilde\ replicate +
group}. For the interaction, use \texttt{\textasciitilde\ knockout * drug}.

%=====================================================================
\section{Chapter 1 summary}
%=====================================================================

Chapter~1 showed what BARCS measures and how it describes an experiment:
\begin{enumerate}
\item It turns read counts into proportions (Step~1).
\item It compares groups with a fold change and its logarithm (Step~2).
\item It works on the log-odds scale (Step~3).
\item It separates counting noise from biological noise (Step~4).
\item It describes the experiment with a design matrix (Step~5).
\end{enumerate}
Chapter~2 shows how BARCS calculates the effect and decides if it is real.

\subsection*{The symbols of Chapter 1}

\begin{center}
\small
\begin{tabular}{@{}l p{6.3cm} l@{}}
\toprule
Symbol & Meaning & Value in our example\\
\midrule
$K$ & read count of the guide in one library & 180, 240, 120, 80\\
$N$ & total read count of one library & 1.0, 1.2, 1.0, 0.8 million\\
$p$ & proportion, $K/N$ & 191 and 111 per million\\
fold change & $p_{\text{drug}} \div p_{\text{control}}$ & 0.58\\
$\log_2$ & number of doublings & $\log_2(0.58) = -0.78$\\
$\ln$ & natural logarithm (powers of $e = 2.718$) & $\ln(0.58) = -0.54$\\
log-odds & $\ln\bigl(p/(1-p)\bigr)$ & $-8.56$ (control)\\
$\mu$ & expected proportion & 0.000191 (control)\\
$\rho$ & biological noise & 0.000000386\\
$X$ & design matrix: a row per library, a column per $\beta$ & 4 rows, 2 columns\\
$\beta_0$ & control log-odds in the model & $-8.565$\\
$\beta_1$ & effect of the drug & $-0.542$\\
$\eta$ & predicted log-odds of a library (row of $X$ times $\beta$) & $-8.565$, $-9.107$\\
\bottomrule
\end{tabular}
\end{center}


\newpage
%=====================================================================
\section{Quiz for Chapter 1}
%=====================================================================

Choose the best answer. The answers and calculations are on the next pages.

\begin{enumerate}[label=\textbf{Q\arabic*.},leftmargin=2.6em,itemsep=6pt]

\item What does one guide do in a CRISPR screen?
\begin{choices}
\item It counts the cells.
\item It adds a drug to the cells.
\item It reads the DNA.
\item It switches off one gene.
\end{choices}

\item A library has $K = 300$ and $N = 1{,}500{,}000$. What is the count per
million?
\begin{choices}
\item 200
\item 300
\item 450
\item 500
\end{choices}

\item Why do we divide each count by its library total $N$?
\begin{choices}
\item To make the numbers bigger.
\item Because the machine reads some libraries more than others.
\item To remove the biological noise.
\item Because BARCS cannot use whole numbers.
\end{choices}

\item A guide has 400 per million in control and 100 per million with a
drug. What is the fold change?
\begin{choices}
\item 4
\item 0.4
\item 0.25
\item 300
\end{choices}

\item What is $\log_2(8)$?
\begin{choices}
\item 2
\item 3
\item 4
\item 8
\end{choices}

\item A gene has $\log_2(\text{fold change}) = -1$. What happened to its
guide?
\begin{choices}
\item It doubled.
\item It did not change.
\item It dropped to half.
\item It dropped to zero.
\end{choices}

\item For a very small proportion $p$, which statement is true?
\begin{choices}
\item The odds are almost equal to $p$.
\item The odds are almost equal to 1.
\item The odds are twice $p$.
\item The odds are always 0.
\end{choices}

\item In the variance formula, what does $\rho = 0$ mean?
\begin{choices}
\item There is no noise at all.
\item There is only counting noise.
\item There is only biological noise.
\item The guide has zero reads.
\end{choices}

\item Use $N\mu(1-\mu)$ with $N = 100$ and $\mu = 0.5$. What is the counting
variance?
\begin{choices}
\item 25
\item 50
\item 0.25
\item 100
\end{choices}

\item The machine reads each library ten times more deeply. What happens to
the biological noise?
\begin{choices}
\item It disappears.
\item It becomes ten times smaller.
\item It does not go away. Only more replicates help.
\item It becomes zero for control guides.
\end{choices}

\item In the design matrix of our example, which $\beta$ belongs to the
column of 1s?
\begin{choices}
\item $\beta_1$
\item the interaction
\item $\beta_0$
\item $\rho$
\end{choices}

\item BARCS gives $\beta_1 = 0.69$ for a guide. Use $e^{0.69} = 2.0$. What
does this mean?
\begin{choices}
\item The guide halves with the drug.
\item The guide does not change.
\item The guide disappears.
\item The guide doubles with the drug.
\end{choices}

\item An experiment has 3 control libraries and 3 drug libraries. The
formula is \texttt{\textasciitilde\ group}. What is the size of the design
matrix?
\begin{choices}
\item 3 rows and 2 columns
\item 6 rows and 2 columns
\item 2 rows and 6 columns
\item 6 rows and 3 columns
\end{choices}

\item A library has the row $(1, 1)$ in the design matrix. The values are
$\beta_0 = -8.0$ and $\beta_1 = -0.5$. What is the predicted log-odds?
\begin{choices}
\item $-8.5$
\item $-8.0$
\item $-7.5$
\item $4.0$
\end{choices}

\item In the knockout-and-drug design, which group has a 1 in the
interaction column?
\begin{choices}
\item normal, no drug
\item knockout, no drug
\item normal, drug
\item knockout, drug
\end{choices}

\item You add a paired-replicate column to the model of our four libraries.
How many degrees of freedom stay?
\begin{choices}
\item 3
\item 2
\item 4
\item 1
\end{choices}

\end{enumerate}

\newpage
%=====================================================================
\section{Answers for Chapter 1}
%=====================================================================

\begin{enumerate}[label=\textbf{Q\arabic*.},leftmargin=2.6em,itemsep=6pt]
\item \textbf{D.} Each guide switches off one gene (Section~1).
\item \textbf{A.} $300 \div 1{,}500{,}000 = 0.0002$. Then
  $0.0002 \times 1{,}000{,}000 = 200$ per million.
\item \textbf{B.} Libraries get different numbers of reads. Division by $N$
  makes them comparable (Step~1).
\item \textbf{C.} $100 \div 400 = 0.25$. The guide keeps one quarter of its
  control level.
\item \textbf{B.} $2 \times 2 \times 2 = 8$. You double three times, so
  $\log_2(8) = 3$.
\item \textbf{C.} $-1$ means one halving (Step~2).
\item \textbf{A.} $1 - p$ is almost 1, so $p/(1-p)$ is almost $p$
  (Step~3).
\item \textbf{B.} The extra factor becomes $1 + (N-1) \times 0 = 1$. Only
  counting noise stays (Step~4).
\item \textbf{A.} $100 \times 0.5 \times (1 - 0.5) = 100 \times 0.5 \times
  0.5 = 25$.
\item \textbf{C.} The factor $1 + (N-1)\rho$ grows with $N$. Deeper reading
  does not remove biological noise (Step~4).
\item \textbf{C.} The column of 1s is the intercept. It belongs to
  $\beta_0$, the starting level (Step~5).
\item \textbf{D.} The fold change is $e^{0.69} = 2.0$. The guide doubles.
\item \textbf{B.} One row for each library ($3 + 3 = 6$). One column for the
  intercept and one for drug (Step~5).
\item \textbf{A.} $1 \times (-8.0) + 1 \times (-0.5) = -8.5$ (Step~5).
\item \textbf{D.} The interaction column is knockout times drug. Only
  $1 \times 1 = 1$ (Step~5).
\item \textbf{D.} The model now has 3 columns. So
  $\text{df} = 4 - 3 = 1$ (Step~5).
\end{enumerate}
